\documentclass[11pt]{article}

\usepackage[a4paper,margin=1in]{geometry}
\usepackage[T1]{fontenc}
\usepackage[utf8]{inputenc}
\usepackage{lmodern}
\usepackage{amsmath,amssymb,amsthm}
\usepackage{hyperref}
\usepackage{enumitem}
\setlength{\emergencystretch}{3em}

\title{The Evolution of the Lean Formalization of Appendix B}
\author{Working draft}
\date{\today}

\begin{document}
\maketitle

\begin{abstract}
This note records the evolution of the Lean formalization of Appendix~B in the
\texttt{PptFactorization} development.  Its purpose is not to reproduce hidden
proof search, but to present a publishable mathematical narrative reconstructed
from visible Lean artifacts: theorem statements, supplier lemmas, frontier
changes, build checks, and axiom audits.  The central theme is that the Lean
development did not progress by a single monolithic proof, but by a sequence of
closure layers that steadily replaced broad conditional interfaces with sharper
and more theorem-facing endpoints.
\end{abstract}

\section{Introduction}

Appendix~B studies concentration and large-deviation bounds for normalized
matrix moments on the Hilbert--Schmidt sphere.  In the paper, the proof mixes
geometric concentration on high-dimensional spheres, deterministic local
expansion estimates, and random-matrix control of background and spike terms.
In Lean, the main challenge is not merely to translate the final estimate into
type-correct syntax.  One must decide how to decompose the argument into
suppliers that are mathematically meaningful, auditable, and stable under
later refinements.

The present draft documents that decomposition.  It is meant to evolve with the
formalization itself: whenever a visible hypothesis disappears from a sharp
endpoint, or is replaced by a smaller and more concrete family of suppliers,
this note should explain what changed mathematically and what Lean now certifies.

\section{Mathematical Target}

The Appendix~B upper-bound argument can be summarized as follows.  Let
\[
  X \in S_{HS}(d^2,s)
\]
be drawn from the uniform spherical law on the Hilbert--Schmidt sphere of
complex matrices, with real ambient dimension
\[
  m = 2 d^2 s.
\]
For a fixed moment order \(k\) and a deviation parameter \(\varepsilon>0\), one
seeks an upper large-deviation estimate for the event that the normalized
background moment deviates above its deterministic center by at least
\(\varepsilon\), at the correct exponential speed.

At a high level, the proof has four ingredients:
\begin{enumerate}[leftmargin=*]
  \item positivity of a concrete witness event contained in the desired upper
  tail;
  \item spherical concentration at the real dimension \(2 d^2 s\);
  \item high-probability control of the background typical set;
  \item deterministic control of mixed local-expansion terms near the sharp
  perturbative radius.
\end{enumerate}

The formalization mirrors this structure, but it does so through successive
interfaces rather than a single theorem.  The live Lean frontier therefore has
to be described theorem by theorem.

\section{Formalization Strategy}

The proof audit separates three layers.

First, there is a \emph{conditional geometric pipeline}.  This is the cleanest
abstract interface: given a spherical isoperimetric input, a measurable
background typical set of half mass, a mixed-remainder bound, and the relevant
scalar limits, the endpoint returns the eventual upper limsup bound.

Second, there is a \emph{concrete theorem-facing layer}.  Here the abstract
sequences \(N_d\), \(s_d\), \(\tau_d\), and the spherical law are instantiated
to the actual matrix model.  This is where concrete target probabilities and
dimension transport enter.  The word ``transport'' is important: a pointwise
statement at size \(d\) may use the varying types \(\mathrm{Fin}(d)\) and
\(\mathrm{Fin}(s(d))\), whereas a fixed-type theorem cannot honestly prove
that a fixed finite dimension is eventually equal to \(d^2\).

Third, there is a \emph{frontier-refinement layer}.  This is the most dynamic
part of the project.  Broad assumptions such as a packaged target witness or a
single mixed-word estimate are replaced by smaller suppliers that match the
mathematics more closely: one-sided positivity of the actual upper event,
branchwise mixed-word bounds, and grouped scalar defect limits.

\section{Evolution of the Upper-Bound Frontier}

\subsection{From abstract closure to concrete endpoints}

The first major milestone was to isolate a stable upper-bound closure file with
project-axiom-free wrapper theorems.  Those wrappers did not prove the hard
random-matrix upper large deviation statement on their own, but they exposed
its genuine inputs cleanly.  In particular, the upper pipeline kept visible the
geometric input, the target-event witness data, the background half-mass
estimate, and the mixed local-expansion control.

This phase matters because it fixed the shape of the proof.  Once the abstract
pipeline was stable, later work could close leaves one by one without changing
the final theorem every round.

\subsection{Casewise mixed-word control}

The next step was to replace broad mixed-remainder assumptions by explicit
branchwise estimates.  Instead of a single black-box statement saying that all
mixed words are small, the Lean file now works with three mathematical cases:
\begin{enumerate}[leftmargin=*]
  \item words with exactly one linear defect;
  \item words with exactly one quadratic defect;
  \item all remaining mixed words with multiple defects.
\end{enumerate}

Lean packages these estimates in two different ways.  One package is
theorem-facing and close to the underlying proof: it keeps the three cases
separate.  Another package is endpoint-facing and convenient for the curated
upper theorem: it folds the three cases into a single mixed-word envelope bound.
The point of the refactor is not cosmetic.  It makes the remaining analytic
obligations explicit in the same form in which a mathematician would naturally
prove them.

\subsection{One-sided positivity of the actual upper event}

Another important refinement was to stop treating the upper target event
through an arbitrary witness set alone.  The development now includes a
one-sided positivity supplier for the actual upper-tail event.  This supplier
can be built either from an eventual positive lower bound on the one-sided
upper tail itself, or from the positive mass of a concrete one-column
favourable event together with a deterministic inclusion theorem.

This is conceptually cleaner.  The formal upper endpoint can now say that its
positivity input is the actual one-sided event from the paper, rather than an
auxiliary witness chosen for convenience.

\subsection{The audit wrapper and the corrected frontier}

The most recent closure layer introduces an audit-facing endpoint in which five
older visible leaves are no longer exposed separately:
\[
  hTarget,\quad hWordBound,\quad hOneLinear,\quad hOneQuadratic,\quad hMulti.
\]
They are replaced, in that wrapper, by the smaller grouped inputs
\[
  hOneSided,\quad hWordCases,\quad hDefectLimits,
\]
together with geometric and moment supplier packets.  This is a regrouping of
the wrapper interface, not a proof that all underlying analytic estimates have
disappeared.

Mathematically, this does not mean that the random-matrix analysis has
disappeared.  It means that the theorem interface now reflects the structure of
the proof more faithfully: actual upper-event positivity, grouped branchwise
mixed-word control, and grouped scalar defect limits.

The latest signature audit also corrected an important possible
misinterpretation.  The theorem
\begin{quote}\small\ttfamily
\detokenize{upper_eventual_from_concrete_sequences_of_fullIso_activeTicketSuppliers}
\end{quote}
is an audit wrapper, not the primitive mathematical frontier.  Its hypotheses
\[
  hOneSided,\quad hMoment,\quad hWordCases,\quad hDefectLimits
\]
are wrapper-level packets: they collect several more elementary supplier
statements into compact arguments.  They should not be counted as four new
independent analytic abysses.  The sharper theorem-facing endpoint is
\begin{quote}\small\ttfamily
\detokenize{upper_eventual_from_concrete_sequences_of_fullIso}\\
\detokenize{_positiveDeviationWitnessData_localExpansionMomentWordEnvelope}\\
\detokenize{_caseMixedWordBounds_caseMixedTermLimits}
\end{quote}
That endpoint exposes the real proof shape: witness positivity for the target
event, half-mass of the background typical set, word-envelope bounds and their
term limits, scalar budget/envelope inequalities, and the three mixed-word plus
three mixed-term branches.

The same audit clarified the dimension hypotheses.  The file contains the
pointwise concrete bridge
\[
\texttt{upperConcreteRealDim\_eq\_concreteModel\_realDim},
\]
which identifies the real dimension of the actual size-\(d\) model with
\[
2\,\mathrm{bipartiteDimension}(\mathrm{Fin}(d),\mathrm{Fin}(d))\,
  \mathrm{sampleDimension}(\mathrm{Fin}(s(d))).
\]
It also proves the no-go statement
\[
\texttt{not\_eventually\_fixed\_bipartiteDimension\_eq\_dimensionSquared}.
\]
Thus a hypothesis asserting that a fixed finite matrix type has bipartite
dimension \(d^2\) eventually is not a hard theorem waiting for analysis; it is
a sign that the endpoint is being read through a fixed-type compatibility
wrapper rather than through the pointwise concrete model.

\section{What Lean Currently Certifies}

At the time of writing, the Lean development contains two especially important
upper endpoints.

The first is a sharp theorem-facing endpoint in which the moment-side mixed
remainder has itself been reduced to word-envelope bounds and term limits.
This endpoint is mathematically close to the deterministic proof and therefore
useful for understanding what remains to be established analytically.

The second is the current audit-wrapper endpoint.  It is more compact and
better suited for readback, but it should not be mistaken for the primitive
frontier.  Its visible inputs include one-sided upper-event positivity, the
full spherical isoperimetric supplier, two fixed-type dimension compatibility
bridges, the concrete moment package, and grouped mixed-word and scalar-defect
packages.  The fixed-type bridges are bookkeeping/routing artefacts; the
pointwise concrete dimension bridge is already present in Lean, while the
fixed-type eventual version is explicitly shown not to be closable as stated.

The targeted upper module build
\[
\texttt{lake build PptFactorization.AppendixBUpperBoundClosure}
\]
passes for this state of the development.

For the public theorems added on this route, the axiom audit reports only the
standard Lean and mathlib foundations:
\[
  \texttt{propext},\qquad
  \texttt{Classical.choice},\qquad
  \texttt{Quot.sound}.
\]
Thus these closure layers are fully formalized in Lean/mathlib with no
project-specific axioms.

\section{Remaining Frontier}

The remaining work splits into two different categories.

The first category is routing and import work.  If the full spherical
isoperimetric theorem has already been formalized elsewhere and audited from
mathlib foundations, the task is to import or adapt that supplier, not to
reprove spherical isoperimetry.  Similarly, the fixed-type dimension hypotheses
should be avoided or discharged by moving through the pointwise concrete model
where the dimensions are definitionally the right ones.

The second category is theorem-strength work.  On the upper side, the remaining
analytic suppliers are the concrete positivity route for the target event when
not supplied by the favourable-event construction, the high-probability
half-mass estimate for the background typical set, and the deterministic
word-envelope and scalar term-limit estimates that feed the mixed local
expansion bound.

This is exactly the kind of frontier one wants after a successful structural
formalization.  The Lean file no longer hides those obligations behind vague
interfaces.  It displays them in theorem signatures that correspond to
recognizable mathematical statements.

\section{Methodological Lesson}

The main methodological lesson is that large analytic formalizations benefit
from a layered endpoint strategy.  A broad conditional theorem is useful early,
because it isolates the architecture of the proof.  Later, however, progress
should be measured by visible frontier changes: which assumptions disappeared,
which ones were replaced by smaller suppliers, and which ones remain genuinely
hard mathematics.

That is the sense in which this article is meant to ``follow'' the proof.  It
does not attempt to reproduce hidden proof search.  Instead, it records the
public mathematical evolution of the Lean development, one auditable closure
step at a time.

\end{document}
